\section{Algorithm}

\subsection{Per-tx state}

Every transaction enters LP-210 carrying:

\begin{table}[h]
\centering
\small
\begin{tabular}{@{}lll@{}}
\toprule
\textbf{Field} & \textbf{Type} & \textbf{Source} \\
\midrule
\texttt{tx\_bytes} & ZAP buffer (pointer-shared) & LP-208 header payload \\
\texttt{tx\_index} & uint32 & LP-209 wave ordering position \\
\texttt{snapshot} & ZapDB MVCC version pointer & LP-202 \texttt{state.NewSnapshot()} \\
\texttt{read\_set} & sorted slice of \texttt{(key, version)} & populated during execute \\
\texttt{write\_set} & sorted slice of \texttt{(key, new\_value, prior\_version)} & populated during execute \\
\texttt{verdict} & enum \{pending, committed, rolled-back, invalid\} & set in commit phase \\
\bottomrule
\end{tabular}
\caption{Per-tx state at LP-210 entry.}
\label{tab:pertx}
\end{table}

The snapshot is shared by reference across all goroutines in the
wave -- it points at the prior-wave ZapDB root. Concurrent reads
against the same snapshot are lock-free (ZapDB MVCC invariant).
Writes happen in per-tx-local MVCC overlays.

\subsection{Phase 1 -- speculative execute}

Dispatch $N$ transactions to a worker pool of
\texttt{min(GOMAXPROCS, $N$)} goroutines. Each goroutine runs:

\begin{verbatim}
func executeTx(tx Tx, snap *ZapDBSnapshot,
               rs *ReadSet, ws *WriteSet) {
    ctx := newTxContext(snap, rs, ws)
    runVMOpcodes(tx.bytes, ctx)
    // ctx.rs and ctx.ws now populated
}
\end{verbatim}

Every state access through \texttt{ctx} records into \texttt{rs}
(for reads) or \texttt{ws} (for writes). Read records the
\texttt{(key, version)} pair from the snapshot; write records
\texttt{(key, new\_value, prior\_version)} where
\texttt{prior\_version} is the version observed at write time.

No locks during speculative execute. The ZAP buffer is immutable;
the ZapDB snapshot is immutable from the goroutine's POV. The only
mutation is into \texttt{rs} and \texttt{ws}, which are per-goroutine
state.

\subsection{Phase 2 -- commit-phase conflict detection}

After all $N$ goroutines join, walk the transactions in LP-209 order
(\texttt{tx\_index} ascending). For each $\textrm{tx}_i$, check:

\begin{equation}
\textrm{conflict}(\textrm{tx}_i, \textrm{prior\_committed}) \iff
\exists\,(k, v) \in \textrm{tx}_i.\textrm{read\_set}.\;
\exists\,\textrm{tx}_j \in \textrm{prior\_committed},\;
j < i,\; k \in \textrm{tx}_j.\textrm{write\_set}
\end{equation}

If no conflict: commit $\textrm{tx}_i$'s \texttt{write\_set} to the
staging MVCC overlay. Mark \texttt{verdict = committed}. If
conflict: $\textrm{tx}_i$'s \texttt{rs} and \texttt{ws} are stale.
Mark \texttt{verdict = rolled-back}. $\textrm{tx}_i$ is re-queued
for phase~3.

The check is $O(|\textrm{tx}_i.\textrm{read\_set}|)$ per tx using a
hashmap of keys $\to$ most recent committed writer. Total
commit-phase cost is
$O(\sum_i |\textrm{tx}_i.\textrm{read\_set}|)$ -- linear in total
reads across the wave.

\subsection{Phase 3 -- re-execute rolled-back txs in dependency order}

Re-queued transactions are re-executed sequentially in
\texttt{tx\_index} order. Each re-execution sees the committed
overlay from Phase~2 (which now includes the writes that caused the
rollback). Re-execution is serial because the contended set is small
(in the measured regulated-partner workload, $\sim 5\%$ of txs per wave);
parallelizing re-execution recovers little throughput at significant
complexity cost.

Re-executed transactions may still fail (a tx that depended on a
now-overwritten balance may now underflow). Failed re-execution
marks \texttt{verdict = invalid} and the tx is dropped from the wave
commit. LP-209's wave ordering is independent of validity -- an
invalid tx is recorded in the wave but does not contribute to the
post-wave ZapDB root.

\subsection{Phase 4 -- atomic commit}

The staging MVCC overlay (the union of all committed
\texttt{write\_set}s) is applied to ZapDB as a single MVCC commit.
This is the LP-202 atomic commit primitive: either the entire wave's
writes land or none do.

\begin{verbatim}
overlay := zapdb.NewOverlay(snap)
for _, tx := range wave.txs {
    if tx.verdict == committed {
        overlay.Apply(tx.write_set)
    }
}
overlay.Commit()  // LP-202 O(1) atomic; new ZapDB root published
\end{verbatim}

LP-217's cert ceremony (one cert per wave per LP-209) covers the
new root pointer. LP-211 (cross-shard) interposes here if the wave
is part of a cross-shard cert.

\subsection{Deterministic output proof}

LP-210 introduces no wire format. Every transaction in the wave is
already a ZAP frame (LP-022). The STM machinery is purely runtime
-- MVCC version pointers are ZapDB-internal; read\_set and
write\_set are runtime in-memory structures, never serialized,
never gossiped, never content-addressed.

This is a load-bearing property: a validator implementing LP-210 in
Go and another implementing it in Rust over the same ZAP buffer and
ZapDB state produce \textbf{byte-identical} post-wave roots. The
order of re-executions in Phase~3 is determined by tx\_index order
(LP-209 canonical); the set of committed writes is determined by
the deterministic conflict-detection rule; the MVCC root is
therefore a pure function of $(\textrm{wave\_bytes},
\textrm{prior\_root})$.
